\chapter{System Context and Safety Distinctions}
\label{ch:ch03_background}

% --- Merged from ch05_orius_system_context.tex ---

\section{ORIUS System Context}
\label{ch:orius-context}

Grid-scale battery energy storage systems (BESS) sit at the intersection of
three economic forces: variable renewable generation that must be absorbed,
load profiles that demand dispatchable power on sub-hourly timescales, and
wholesale markets that price energy at volatile clearing rates.  This chapter
positions the ORIUS framework within that operational context, traces the
closed decision loop from forecast to measurement, and identifies the
\emph{Optimize--Dispatch boundary} as the location where observation quality
most directly threatens safety.

% ----------------------------------------------------------------
\subsection{Why Grid-Scale Batteries Matter}

The penetration of wind and solar generation has transformed balancing-area
operations.  In 2024 the ERCOT interconnection recorded negative real-time
prices for 18\% of daytime hours, while evening ramp rates exceeded
12\,GW\,h$^{-1}$ on 37 days.  Grid-scale batteries absorb these mismatches:
they charge when marginal cost is low, discharge when it is high, and provide
ancillary services (frequency regulation, spinning reserve) that stabilise the
grid during transients.

Three properties distinguish BESS from conventional peakers:
\begin{enumerate}
  \item \textbf{Bidirectional power flow.}  A 100\,MW / 400\,MWh installation
    can charge or discharge at full rate within seconds, making it a
    controllable load \emph{and} a controllable generator.
  \item \textbf{State-of-charge coupling.}  Every dispatch decision at time~$t$
    constrains the feasible set at time~$t+1$ through the SOC dynamics.  The
    controller cannot treat time steps independently.
  \item \textbf{Degradation risk.}  Deep cycling, high C-rates, and
    over-voltage operation accelerate cell ageing.  Safety envelopes are
    therefore both \emph{hard} (BMS trip thresholds) and \emph{soft}
    (warranty-preserving SOC margins).
\end{enumerate}

These properties mean that dispatch quality depends not only on forecast
accuracy but on the fidelity of the real-time observation that informs each
control step.

% ----------------------------------------------------------------
\subsection{ORIUS as a Closed Decision Loop}

The ORIUS system implements a five-stage closed loop:

\[
  \underbrace{\text{Forecast}}_{\hat{y}_{t+1:t+H}}
  \;\longrightarrow\;
  \underbrace{\text{Optimize}}_{a_t^\star}
  \;\longrightarrow\;
  \underbrace{\text{Dispatch}}_{a_t^{\mathrm{exec}}}
  \;\longrightarrow\;
  \underbrace{\text{Measure}}_{\tilde{z}_t}
  \;\longrightarrow\;
  \underbrace{\text{Monitor}}_{\mathcal{C}_t}
  \;\longrightarrow\;
  \text{(next cycle)}
\]

\begin{description}
  \item[Forecast.]  A model family $\mathcal{F}$ produces point forecasts and
    prediction intervals for load, wind, solar, and price over a rolling
    horizon $H$ (typically 24--48\,h at hourly resolution).

  \item[Optimize.]  A dispatch optimiser receives the forecast distribution
    and the current SOC estimate, then solves a constrained programme to
    produce a candidate action $a_t^\star = (\text{charge}_t,
    \text{discharge}_t)$.

  \item[Dispatch.]  The candidate action is checked against safety constraints
    (SOC bounds, ramp limits, BMS thresholds).  If the action is feasible it
    is forwarded to the battery management system; otherwise it is repaired to
    the nearest safe action $a_t^{\mathrm{safe}}$.

  \item[Measure.]  Operational telemetry $\tilde{z}_t$ returns from the plant:
    actual SOC, terminal voltage, cell temperatures, grid-side power meter
    readings.  This measurement may be delayed, dropped, corrupted, or stale.

  \item[Monitor.]  A certificate $\mathcal{C}_t$ records the action taken, the
    observation quality score, the interval width, and whether the step passed
    all safety invariants.  The certificate feeds the next forecast cycle and
    is archived for audit.
\end{description}

% ----------------------------------------------------------------
\subsection{The Optimize--Dispatch Boundary as the True Risk Location}

Not all stages of the loop carry equal risk.  A forecast error at
$t + 12$\,h can be corrected by re-optimisation at $t + 1$; a measurement
delay at the \emph{Measure} stage is visible to the \emph{Monitor} and can
trigger a conservative hold.  The highest-risk transition is the
\textbf{Optimize $\to$ Dispatch} boundary, for three reasons:

\begin{enumerate}
  \item \textbf{Commitment.}  Once the BMS receives a power setpoint, the
    battery executes it within milliseconds.  The action is irreversible on the
    control-cycle timescale.
  \item \textbf{Observation dependence.}  The optimiser's candidate~$a_t^\star$
    is conditioned on the most recent observation~$\tilde{z}_{t-1}$.  If that
    observation is stale or corrupted, the candidate may lie outside the true
    feasible set without any visible warning.
  \item \textbf{Constraint coupling.}  An unsafe action at time~$t$ shifts the
    SOC into a region that may render \emph{all} future actions infeasible
    until a recovery period elapses (the ``SOC trap'').
\end{enumerate}

The DC3S adapter (Chapter~\ref{ch:dc3s-adapter}) is positioned precisely at
this boundary: it intercepts~$a_t^\star$, inflates the uncertainty set to
account for observation degradation, and projects the action into the
tightened safe set before it reaches the BMS.

% ----------------------------------------------------------------
\subsection{Operational Telemetry Pathways}

Telemetry in a grid-scale BESS traverses multiple communication layers before
reaching the dispatch controller:

\begin{enumerate}
  \item \textbf{Cell-level sensors} (voltage, temperature, current) report to
    the Battery Management System (BMS) over CAN bus at 1--10\,Hz.
  \item \textbf{BMS aggregation} computes pack-level SOC, state-of-health
    (SOH), and thermal limits, then publishes to the site SCADA at 1\,Hz.
  \item \textbf{SCADA gateway} relays readings to the Energy Management System
    (EMS) via Modbus/TCP or IEC~61850.  This link is typically on-site LAN
    with sub-second latency.
  \item \textbf{Cloud telemetry} streams from the EMS to the dispatch
    controller (ORIUS) over cellular, fibre, or satellite uplink.  This
    link is the primary source of packet loss, delay jitter, and reordering.
  \item \textbf{Grid-side meters} (revenue meters, PMUs) provide independent
    power and frequency measurements at 4\,s or sub-second resolution.
\end{enumerate}

Each hop introduces a distinct failure mode: CAN bus saturation at high
C-rates, SCADA gateway reboot during firmware updates, cellular dropout during
storms, and satellite latency during congestion.  The telemetry fault taxonomy
used in CPSBench (Chapter~\ref{ch:cpsbench}) mirrors these real failure modes.

% ----------------------------------------------------------------
\subsection{Why Measurement Quality Enters the Control Problem}

Classical dispatch formulations assume that the controller observes the true
state~$z_t$ at each step.  In practice the controller observes a degraded
version~$\tilde{z}_t$ that may differ from~$z_t$ in timing, magnitude, or
completeness.  The gap between $z_t$ and $\tilde{z}_t$ has three consequences
for safety:

\begin{enumerate}
  \item \textbf{SOC estimation error.}  If the reported SOC is stale by
    $\Delta t$ steps, the true SOC has drifted by up to
    $P_{\max} \cdot \Delta t / E_{\mathrm{cap}}$ in normalised units, where
    $P_{\max}$ is the rated power and $E_{\mathrm{cap}}$ the nameplate energy.
  \item \textbf{Forecast conditioning error.}  The most recent load and
    renewables readings condition the next forecast window.  Corrupted readings
    propagate into biased intervals, undermining conformal coverage guarantees.
  \item \textbf{Constraint tightening mismatch.}  A robust optimiser that
    tightens constraints based on \emph{assumed} interval width may
    under-tighten if the true observation quality is worse than assumed, or
    over-tighten if quality is better, sacrificing economic value.
\end{enumerate}

The central thesis of this dissertation is that measurement quality must be
\emph{explicitly scored and fed back} into the uncertainty set at each control
step.  The observation quality estimator (OQE, Chapter~\ref{ch:dc3s-adapter})
and the reliability-adaptive conformal certificate (RAC-Cert) implement this
feedback loop.

% ----------------------------------------------------------------
\subsection{Chapter Summary}

This chapter established the ORIUS system as a five-stage closed decision
loop---Forecast, Optimize, Dispatch, Measure, Monitor---operating over
grid-scale battery storage.  The Optimize--Dispatch boundary was identified as
the highest-risk transition because actions are irreversible, observation-dependent,
and coupled through SOC dynamics.  Telemetry traverses multiple communication
layers, each introducing distinct degradation modes.  Because measurement
quality directly affects SOC estimation, forecast conditioning, and constraint
tightening, it must enter the control formulation as a first-class variable.
The remainder of Part~II develops the mathematical machinery to accomplish this.


% --- Merged from ch06_data_telemetry_scope.tex ---

\section{Data, Telemetry, and Dataset Scope}
\label{ch:data-telemetry}

Every empirical claim in this dissertation rests on two locked dataset
surfaces: the Germany OPSD archive and the United States EIA-930 feed.  This
chapter documents the signals, their provenance, the telemetry pathways
through which they enter the ORIUS loop, and the fault modes that make
degraded observation an ordinary operating condition rather than an edge case.

% ----------------------------------------------------------------
\subsection{Germany and United States Dataset Surface}

The ORIUS evaluation spans two geographically and regulatorily distinct
power systems.

\paragraph{Germany (OPSD).}
The Open Power System Data platform aggregates hourly generation, load, and
day-ahead price series from ENTSO-E and the German TSOs.  The locked dataset
covers 2018-10-07 to 2020-09-30 (17\,377 hourly rows), with four primary
targets: \texttt{load\_mw}, \texttt{wind\_mw}, \texttt{solar\_mw}, and
\texttt{price\_eur\_mwh}.  Weather covariates (temperature, irradiance, wind
speed) are sourced from the Open-Meteo archive API and joined on UTC
timestamp.  The full feature matrix contains 94 engineered columns after
calendar encoding and lag construction.

\paragraph{United States (EIA-930).}
The U.S.\ Energy Information Administration Form~930 provides hourly
demand and net generation by source for each balancing authority.  Three
regions are retained:

\begin{table}[ht]
\centering
\caption{Dataset cards for the locked evaluation surfaces.  Source:
\texttt{reports/publication/dataset\_cards.csv}.}
\label{tab:dataset-cards}
\begin{tabular}{llrrrl}
\toprule
\textbf{Key} & \textbf{Region} & \textbf{Rows} & \textbf{Features} &
  \textbf{Coverage} & \textbf{Period} \\
\midrule
DE       & Germany (OPSD)    & 17\,377 & 94  & 100\% & 2018-10 to 2020-09 \\
US\_ERCOT & USA (EIA-930)    & 13\,453 & 62  & 100\% & 2024-07 to 2026-01 \\
US\_MISO  & USA (EIA-930)    & 13\,638 & 114 & 100\% & 2024-07 to 2026-01 \\
US\_PJM   & USA (EIA-930)    & 13\,639 & 62  & 100\% & 2024-07 to 2026-01 \\
\bottomrule
\end{tabular}
\end{table}

Weather covariates for US regions are likewise obtained from Open-Meteo.
Carbon intensity is estimated from the SMARD generation mix (DE) or a
proxy generation-mix model (US).

Complete source Data Cards for the battery witness rows and the broader
ORIUS portability package are provided in Appendix~\ref{app:dataset-cards}.

% ----------------------------------------------------------------
\subsection{Forecast Targets and Operational Channels}

Each dataset surface exposes three operational forecast targets:

\begin{description}
  \item[\texttt{load\_mw}] Net system load (total demand minus behind-the-meter
    solar), measured at the transmission level.  This is the primary signal for
    arbitrage scheduling.
  \item[\texttt{wind\_mw}] Aggregate onshore and offshore wind generation.
    Wind exhibits high variance and heavy tails, particularly during storm
    fronts, making calibration challenging.
  \item[\texttt{solar\_mw}] Aggregate photovoltaic generation.  Solar follows a
    deterministic diurnal envelope modulated by cloud cover; residual
    uncertainty is sharply bimodal (clear vs.\ overcast).
\end{description}

A fourth target, \texttt{price\_eur\_mwh} (DE only), drives the economic
objective in the dispatch optimiser.  For US regions, a simplified tariff
schedule or locational marginal price proxy is used.

All targets are sampled at hourly resolution and aligned to UTC.  The forecast
horizon is $H = 48$\,h with a stride of 1\,h, yielding a rolling-origin
evaluation with 168 test windows per week.

% ----------------------------------------------------------------
\subsection{Observed Telemetry Versus Truth Trajectory}

A central design choice in ORIUS is the separation of \emph{truth state}
and \emph{observed state}.  At every control step~$t$:

\begin{itemize}
  \item The \textbf{truth trajectory} $\{z_t\}$ records what actually happened
    at the plant: the true SOC, the true net load, the true renewable
    generation.  This trajectory is never available to the controller in real
    time; it is logged by the CPSBench harness for post-hoc evaluation.
  \item The \textbf{observed trajectory} $\{\tilde{z}_t\}$ is what the
    controller actually sees after telemetry faults have been applied.  It may
    contain dropped readings (NaN), delayed values, sensor spikes, stale
    repetitions, or drifted baselines.
\end{itemize}

The gap $\|z_t - \tilde{z}_t\|$ is the observation error.  All safety metrics
(TSVR, severity) are computed on the truth trajectory, while the controller's
decisions are computed on the observed trajectory.  This dual-trajectory
architecture is the foundation of the CPSBench evaluation protocol
(Chapter~\ref{ch:cpsbench}).

% ----------------------------------------------------------------
\subsection{Missingness, Corruption, and Signal Degradation}

Telemetry faults fall into six categories, each modelled by a corresponding
CPSBench fault injector:

\begin{table}[ht]
\centering
\caption{Telemetry fault taxonomy and operational causes.}
\label{tab:fault-taxonomy}
\begin{tabular}{lll}
\toprule
\textbf{Fault type} & \textbf{Mechanism} & \textbf{Field cause} \\
\midrule
Dropout       & Packet loss (NaN)       & Cellular dropout, gateway reboot \\
Delay jitter  & Arrival-time shift      & Network congestion, buffering \\
Spike         & Multiplicative outlier  & Sensor transient, ADC overflow \\
Stale sensor  & Frozen reading          & Firmware hang, watchdog failure \\
Drift         & Slow baseline shift     & Sensor ageing, calibration decay \\
Reorder       & Out-of-sequence packets & UDP path diversity, queue inversion \\
\bottomrule
\end{tabular}
\end{table}

These fault types are not mutually exclusive.  The \texttt{drift\_combo}
scenario in CPSBench applies all six simultaneously at configurable severity,
representing a worst-case but operationally plausible field condition.

% ----------------------------------------------------------------
\subsection{Why Telemetry Faults Are Ordinary, Not Edge Cases}

Field data from commercial BESS installations suggests that telemetry
degradation is the norm rather than the exception:

\begin{itemize}
  \item Cellular uplinks in rural wind-farm sites experience 2--8\% packet
    loss during normal operation, rising to 15--25\% during severe weather.
  \item SCADA gateway reboots occur 1--3 times per month, each causing a
    30--90\,s telemetry blackout.
  \item Sensor drift in current transducers accumulates at 0.1--0.5\% of
    full-scale per year, enough to bias SOC estimates by 2--5\,MWh on a
    100\,MW/400\,MWh system within six months.
  \item Firmware updates on BMS controllers occasionally freeze individual
    cell-group readings for minutes until the watchdog timer resets the
    reporting thread.
\end{itemize}

A safety framework that assumes clean telemetry is therefore unsafe by
construction.  The ORIUS evaluation treats faulted telemetry as the
\emph{default} operating condition and clean telemetry as a special case
(the \texttt{nominal} scenario with zero fault injection).

% ----------------------------------------------------------------
\subsection{Artifact Contracts and Profile Lock}

Each dataset surface is frozen at a specific commit hash and validated by a
reproducibility contract:

\begin{enumerate}
  \item \textbf{Row-count invariant.}  The ingestion pipeline asserts that
    the row count matches the locked value in
    \texttt{reports/publication/dataset\_cards.csv}.
  \item \textbf{Column schema.}  The feature matrix schema (column names,
    dtypes, nullable flags) is checked against a frozen JSON schema at
    pipeline entry.
  \item \textbf{Target coverage.}  All three forecast targets must have
    100\% non-null coverage in the locked period.  Any gap causes the pipeline
    to abort with a traceable error.
  \item \textbf{Hash verification.}  The SHA-256 hash of the serialised
    Parquet file is compared to the locked value before any model training or
    evaluation proceeds.
\end{enumerate}

These contracts ensure that no drift in upstream data sources can silently
invalidate the empirical results.

% ----------------------------------------------------------------
\subsection{Scope Boundary}

The dataset scope of this dissertation is intentionally bounded:

\begin{itemize}
  \item \textbf{Included:} hourly bulk-power system signals (load, wind,
    solar, price) for four regions across two continents.
  \item \textbf{Excluded:} sub-hourly (5-min, 1-min) data, distribution-level
    signals, behind-the-meter solar, demand response, and frequency regulation
    signals.
  \item \textbf{Rationale:} the hourly resolution matches the natural
    arbitrage cycle of 2--4\,h duration storage and aligns with the temporal
    granularity of conformal calibration windows.
\end{itemize}

Extending the framework to sub-hourly resolution is discussed in
Chapter~\ref{ch:orius-bench-battery} as part of the benchmark roadmap.

% ----------------------------------------------------------------
\subsection{Chapter Summary}

This chapter documented the two dataset surfaces (Germany OPSD and US EIA-930)
that underpin every empirical result, the three forecast targets (\texttt{load\_mw},
\texttt{wind\_mw}, \texttt{solar\_mw}), and the six telemetry fault types that
degrade observations between the plant and the controller.  The dual-trajectory
architecture---truth state versus observed state---was introduced as the
foundation for safety evaluation.  Artifact contracts and hash-locked profiles
ensure reproducibility.  Telemetry faults are treated as ordinary operating
conditions, not edge cases, motivating the observation-aware machinery
developed in the following chapters.


% --- Merged from ch04_observed_state_safety_illusion.tex ---

\section{The Observed-State Safety Illusion}
\label{ch:safety-illusion}

A battery controller that checks constraints only against its own
observations will, by construction, never detect a violation it cannot
see.  This chapter formalises the \emph{observed-state safety illusion},
shows how it arises in the ORIUS dispatch loop, explains why neither
better forecasting nor tighter optimisation can resolve it, and introduces
the dual-trajectory evaluation paradigm that makes the illusion visible.

\subsection{From Partial Observability to Safety Illusion}

Partial observability is a well-studied concept in control theory and
reinforcement learning.  A partially observable Markov decision process
(POMDP) acknowledges that the agent sees a noisy or incomplete projection
of the true state and plans accordingly, typically via belief-state
updates.  The safety illusion arises when the agent \emph{does not plan
accordingly}---when the controller treats its observation as ground truth
and evaluates constraints against it without accounting for observation
quality.

In a POMDP with a well-calibrated belief, the agent maintains a
distribution over true states and can compute the worst-case constraint
satisfaction probability.  In practice, battery dispatch controllers do
not maintain belief distributions; they receive a point estimate
$\hat{x}_t$ from the SCADA system and optimise against it.  The gap
between the POMDP ideal and the deployed reality is the source of the
illusion.


\subsection{Why the Illusion Arises in the ORIUS Loop}

\subsubsection{The Optimise-to-Dispatch Boundary}

As described in Chapter~\ref{ch:problem-formulation}, the ORIUS loop
has a clear boundary between the optimiser and the dispatcher.  The
optimiser solves~\eqref{eq:dispatch-obj} on the observed state and
produces a candidate action~$a_t^{\star}$.  The dispatcher (in the
baseline case) forwards this action to the plant without modification.

Both the optimiser and the baseline dispatcher operate on
$\hat{x}_t$.  The safety constraints
\[
  \mathrm{SOC}_{\min} \;\le\;
  \widehat{\mathrm{SOC}}_t + \frac{\eta_c [a_t]^{+} - [a_t]^{-}/\eta_d}{E_{\max}} \cdot \Delta t
  \;\le\; \mathrm{SOC}_{\max}
\]
are checked on the observed SOC.  If $\widehat{\mathrm{SOC}}_t$ over- or
under-estimates the true SOC, the constraint check is meaningless with
respect to physical safety.

\subsubsection{Why Good Forecasting Does Not Solve It}

A natural objection is: ``Improve the forecast, and the observation error
disappears.''  This is incorrect for a fundamental reason.  The forecast
model predicts future load, solar, and wind \emph{given current
observations}.  If the current observation is itself degraded---a stale
SOC reading, a dropped load packet---the forecast inherits the error.  A
perfect forecast model with imperfect inputs produces imperfect outputs.

More precisely, the forecast model learns $\mathbb{E}[Y_{t+h} \mid
\hat{X}_t]$, not $\mathbb{E}[Y_{t+h} \mid X_t]$.  When
$\hat{X}_t \ne X_t$, the conditional expectation is evaluated at the wrong
conditioning point.  The resulting action is optimal for a state the plant
does not occupy.

\subsection{Formal Statement of the Illusion}
\label{sec:formal-illusion}

\begin{definition}[Observed-State Safety Illusion]
\label{def:safety-illusion}
Let $\hat{x}_t$ be the observed state, $x_t$ the true state, and
$a_t$ the dispatched action.  The \emph{observed-state safety illusion}
at time~$t$ is the conjunction of three conditions:
\begin{enumerate}
  \item \textbf{Observed safety:}
    $a_t \in \mathcal{A}(\hat{x}_t)$---the action satisfies all
    constraints when evaluated on the observation.
  \item \textbf{True-state violation:}
    $a_t \notin \mathcal{A}(x_t)$---the action violates at least one
    constraint when evaluated on the truth.
  \item \textbf{Undetectability:}
    The controller has no mechanism to detect condition~(2) from its
    available information $\hat{x}_t$ alone.
\end{enumerate}
\end{definition}

Condition~(3) distinguishes the illusion from a detectable fault.  If the
controller could infer that $\hat{x}_t$ is unreliable (e.g., via an OQE
score), it could compensate.  The illusion is \emph{silent}: the
controller proceeds with full confidence in a physically unsafe action.

\begin{theorem}[Existence of the illusion under dropout]
\label{thm:illusion-existence}
Under the fault model of Section~\ref{sec:fault-model} with dropout
rate $\rho > 0$ and last-value hold imputation, there exists a time
step~$t$ and a deterministic LP dispatch action~$a_t$ such that the
observed-state safety illusion (Definition~\ref{def:safety-illusion})
holds with positive probability.
\end{theorem}

\begin{proof}[Proof sketch]
Consider a discharge phase where
$\mathrm{SOC}_t^{\mathrm{true}} = \mathrm{SOC}_{\min} + \epsilon$ for
small $\epsilon > 0$.  If the most recent $k$ SOC readings were dropped
and the last-value hold reports
$\widehat{\mathrm{SOC}}_t = \mathrm{SOC}_{\min} + \epsilon + k \cdot
\delta$ (the SOC at the time of the last successful reading, before $k$
discharge steps each consuming $\delta$), then the LP observes a
comfortable margin and dispatches a further discharge~$a_t < 0$.  The
true SOC after this action is
$\mathrm{SOC}_{t+1}^{\mathrm{true}} = \mathrm{SOC}_{\min} + \epsilon -
|a_t| \cdot \Delta t / (E_{\max} \cdot \eta_d) < \mathrm{SOC}_{\min}$
for $|a_t|$ exceeding $\epsilon \cdot E_{\max} \cdot \eta_d / \Delta t$.
The action satisfies observed constraints, violates true constraints, and
the LP cannot detect the violation---all three conditions hold.
\end{proof}

\subsection{Battery-Domain Realisation Under Telemetry Faults}
\label{sec:battery-narrative}

\subsubsection{A Simple Battery Narrative}

Consider a 100~MWh / 50~MW battery operating over 168~hours (one week)
under the \texttt{dropout\_only} scenario at $\rho = 0.20$.

\begin{enumerate}
  \item \textbf{Hour 0--24.}  Telemetry is clean.  The deterministic LP
    dispatches charge during solar surplus and discharge during evening
    peak.  Observed and true SOC tracks are indistinguishable.
  \item \textbf{Hour 25.}  A dropout burst begins: five consecutive SOC
    readings are missing.  The last-value hold freezes
    $\widehat{\mathrm{SOC}}$ at 55\%.
  \item \textbf{Hour 26--28.}  The LP continues discharging against the
    frozen observation.  True SOC falls from 55\% to 38\% to 24\%.  The
    observed SOC remains at 55\%.
  \item \textbf{Hour 29.}  The LP dispatches a 40~MW discharge, believing
    SOC is 55\%.  True SOC drops to 16\%, breaching the 20\% floor.
    Observed SOC still shows 55\%.  The illusion is complete.
  \item \textbf{Hour 30.}  Telemetry recovers.  The LP observes SOC at
    16\% and abruptly curtails dispatch, but the damage---an SOC
    excursion below the safety floor---has already occurred.
\end{enumerate}

This narrative is not hypothetical.  The CPSBench harness reproduces it
deterministically (seed 42, \texttt{dropout\_only}, $\rho = 0.20$), and
the true-state violation is recorded in the TSVR metric.

\subsection{Why Standard Evaluation Misses the Problem}
\label{sec:standard-eval-fails}

Standard battery dispatch evaluation protocols report:

\begin{itemize}
  \item Revenue or regret (economic performance).
  \item Observed SOC constraint violations (observed safety).
  \item Forecast accuracy (MAPE, RMSE, CRPS).
\end{itemize}

None of these metrics involve the \emph{true} SOC trajectory.  A
controller that achieves zero observed violations can simultaneously have
a high true-state violation rate.  The evaluation is complicit in the
illusion: it checks safety on the same degraded observation that caused
the unsafe action.

\begin{proposition}[Insufficiency of observed-state evaluation]
\label{prop:obs-eval-insufficient}
For any $\mathrm{TSVR} > 0$, there exists a controller-scenario pair
with $\mathrm{OSVR} = 0$ and the given $\mathrm{TSVR}$.  Observed-state
evaluation is therefore \emph{insufficient} to certify physical safety.
\end{proposition}

\begin{proof}
Take any controller $\pi$ and any fault scenario with dropout rate
$\rho > 0$.  Since $\pi$ evaluates constraints on $\hat{x}_t$ and the
telemetry operator preserves observed-state feasibility by construction
(the hold value is always within bounds), $\mathrm{OSVR} = 0$.  By
Theorem~\ref{thm:illusion-existence}, there exists a step where $a_t$
violates true-state constraints, so $\mathrm{TSVR} > 0$.
\end{proof}

\subsection{CPSBench and Dual-Trajectory Evaluation}
\label{sec:cpsbench-dual}

To break the illusion, we introduce the \emph{dual-trajectory} evaluation
paradigm implemented in the CPSBench harness.  At each step~$t$, the
harness maintains two parallel state machines:

\begin{enumerate}
  \item \textbf{Observed trajectory.}  The SOC evolution computed from
    the action~$a_t$ and the observed state~$\hat{x}_t$, using the
    dynamics~\eqref{eq:soc-dynamics} with $\widehat{\mathrm{SOC}}_t$ as
    the initial condition.
  \item \textbf{True trajectory.}  The SOC evolution computed from the
    \emph{same} action~$a_t$ but with the true
    state~$\mathrm{SOC}_t^{\mathrm{true}}$ as the initial condition.
\end{enumerate}

The true trajectory is never visible to the controller; it is maintained
solely for evaluation.  This separation is critical: the controller
cannot ``cheat'' by observing the truth, and the evaluator cannot be
fooled by the observation.

\begin{definition}[True-State Violation Rate (TSVR)]
\label{def:tsvr}
\[
  \mathrm{TSVR} \;=\; \frac{1}{T} \sum_{t=1}^{T}
  \mathbf{1}\!\Bigl[\,
    \mathrm{SOC}_{t}^{\mathrm{true}} < \mathrm{SOC}_{\min}
    \;\;\text{or}\;\;
    \mathrm{SOC}_{t}^{\mathrm{true}} > \mathrm{SOC}_{\max}
  \,\Bigr].
\]
\end{definition}

\begin{definition}[Intervention Rate (IR)]
\label{def:ir}
\[
  \mathrm{IR} \;=\; \frac{1}{T} \sum_{t=1}^{T}
  \mathbf{1}\!\bigl[\,a_t^{\mathrm{safe}} \ne a_t^{\star}\,\bigr],
\]
where $a_t^{\star}$ is the optimiser's proposed action and
$a_t^{\mathrm{safe}}$ is the action actually dispatched after the
DC3S shield.
\end{definition}

The dual-trajectory architecture enables unambiguous attribution: every
true-state violation is traceable to a specific action, a specific
observation error, and a specific gap between $\hat{x}_t$ and $x_t$.

\subsection{Why the Illusion Demands a Runtime Shield}
\label{sec:shield-motivation}

The illusion cannot be resolved at the optimiser level for two reasons:

\begin{enumerate}
  \item \textbf{Information asymmetry.}  The optimiser has access only to
    $\hat{x}_t$.  Without a reliability signal, it cannot distinguish
    clean telemetry from degraded telemetry.  Any fixed margin it
    applies will be either too conservative (wasting revenue when
    telemetry is clean) or too aggressive (missing violations when
    telemetry degrades).
  \item \textbf{Separation of concerns.}  The optimiser's objective is
    economic; the safety layer's objective is physical.  Mixing the two
    creates Pareto conflicts that are better resolved by a dedicated
    post-optimiser shield.  The shield can adapt its tightening in
    real time based on the OQE score, something a batch optimiser
    solved minutes ago cannot do.
\end{enumerate}

The DC3S shield resolves the illusion by introducing an
\emph{observation-quality-aware} constraint tightening that operates at
the dispatch boundary.  The shield does not replace the optimiser; it
corrects the optimiser's blind spot.

\begin{theorem}[DC3S Feasibility Guarantee]
\label{thm:dc3s-informal}
Let $\alpha \in (0,1)$ be the conformal miscoverage level, $n$ the
calibration-set size, and $\bar{w} = T^{-1}\sum_{t=1}^T w_t$ the
episode-mean reliability score.  Under Assumptions~A1, A2, A3, and~A5,
the DC3S shield with conformal quantile $\hat{q}_\alpha$ inflated by
$1/w_t$ satisfies, for any finite episode of $T$ steps:
\[
  \mathrm{TSVR} \;\le\; \alpha\,(1 - \bar{w}) \;+\; \varepsilon_n,
\]
where $\varepsilon_n = O(1/\sqrt{n})$ is the finite-sample calibration
correction term.  In particular:
\begin{enumerate}
  \item At $\bar{w} = 1$ (perfect telemetry) and large $n$:
    $\mathrm{TSVR} \to 0$.
  \item At $\bar{w} < 1$ (degraded telemetry): the bound is
    proportional to the mean reliability gap $(1 - \bar{w})$.
  \item As $n \to \infty$: $\varepsilon_n \to 0$ and the bound
    converges to the episode-mean reliability term exactly.
\end{enumerate}
This result is the single-episode analogue of Theorem~T3
(ORIUS Core Bound, Chapter~\ref{ch:core-bound}), which proves
$\mathbb{E}[V] \le \alpha(1-\bar{w})T$ in expectation over episodes.
\end{theorem}

The formal version of this theorem, with precise assumptions and a
complete proof, is developed in subsequent chapters.

\subsection{Chapter Summary}

The observed-state safety illusion is the mechanism by which the OASG
operates in practice: a controller evaluates constraints on a degraded
observation, concludes the system is safe, and dispatches an action that
is physically unsafe.  The illusion is invisible to any evaluation
protocol that checks constraints on the observed trajectory only.  The
dual-trajectory architecture of CPSBench makes the illusion visible, and
the DC3S runtime shield makes it preventable.  The remainder of the
dissertation develops the shield, proves its guarantees, and validates
them empirically across US and German grid data under the full fault model
of Chapter~\ref{ch:problem-formulation}.


% --- Merged from ch07_battery_dynamics_dispatch.tex ---

\section{Battery Dynamics and Dispatch Problem}
\label{ch:battery-dynamics}

This chapter formalises the state-space model of a grid-scale battery energy
storage system (BESS), derives the dispatch optimisation programme, and
shows how observation uncertainty enlarges the feasible set in ways that
classical deterministic formulations cannot capture.

% ----------------------------------------------------------------
\subsection{Battery State Variables}

The BESS plant is characterised by the following state and parameter
variables at each hourly time step~$t$:

\begin{table}[ht]
\centering
\caption{Battery state and parameter variables.}
\label{tab:battery-vars}
\begin{tabular}{cll}
\toprule
\textbf{Symbol} & \textbf{Name} & \textbf{Units / Range} \\
\midrule
$E_t$          & State of charge (energy)      & MWh, $[0, E_{\max}]$ \\
$P_t^{+}$     & Charge power                  & MW,  $[0, P_{\max}^{+}]$ \\
$P_t^{-}$     & Discharge power               & MW,  $[0, P_{\max}^{-}]$ \\
$\eta^{+}$    & Charge efficiency              & dimensionless, $(0,1]$ \\
$\eta^{-}$    & Discharge efficiency           & dimensionless, $(0,1]$ \\
$E_{\max}$    & Nameplate energy capacity      & MWh \\
$P_{\max}$    & Rated power (symmetric)        & MW \\
$R_{\max}$    & Maximum ramp rate              & MW\,h$^{-1}$ \\
$\mathrm{SOC}_t$ & Normalised SOC, $E_t / E_{\max}$ & $[0, 1]$ \\
\bottomrule
\end{tabular}
\end{table}

For the ORIUS reference plant the nameplate parameters are
$E_{\max} = 10\,000$\,MWh, $P_{\max} = 200$\,MW, $\eta^{+} = \eta^{-} = 0.95$,
and $R_{\max} = P_{\max}$ (no additional ramp constraint beyond rated power).
Soft SOC margins are enforced at $[\mathrm{SOC}_{\min}, \mathrm{SOC}_{\max}]
= [0.05, 0.95]$; hard BMS trips occur at $[0.0, 1.0]$.

% ----------------------------------------------------------------
\subsection{SOC Dynamics}
\label{sec:soc-dynamics}

The energy balance over one time step ($\Delta t = 1$\,h) is:
\begin{equation}
\label{eq:soc-dynamics-battery}
  E_{t+1} = E_t + \eta^{+}\,P_t^{+}\,\Delta t
                 - \frac{P_t^{-}}{\eta^{-}}\,\Delta t.
\end{equation}
In normalised form:
\begin{equation}
\label{eq:soc-normalised}
  \mathrm{SOC}_{t+1} = \mathrm{SOC}_t
    + \frac{\eta^{+}\,P_t^{+} - P_t^{-}/\eta^{-}}{E_{\max}}\,\Delta t.
\end{equation}
This is a deterministic, linear transition---the stochasticity enters through
the \emph{observation} of $E_t$ (which may be stale or corrupted) and through
the forecast of external signals that drive the optimiser's choice of
$(P_t^{+}, P_t^{-})$.

% ----------------------------------------------------------------
\subsection{Power, Ramp-Rate, and SOC Constraints}
\label{sec:constraints}

The feasible action set at time~$t$ is defined by the following constraints:

\paragraph{Power bounds.}
\begin{equation}
\label{eq:power-bounds}
  0 \le P_t^{+} \le P_{\max}^{+}, \qquad
  0 \le P_t^{-} \le P_{\max}^{-}.
\end{equation}

\paragraph{Mutual exclusion.}
Simultaneous charge and discharge is prohibited:
\begin{equation}
\label{eq:mutual-exclusion}
  P_t^{+} \cdot P_t^{-} = 0.
\end{equation}
In the LP relaxation this is enforced via a binary indicator or by solving
separate charge and discharge sub-problems and selecting the dominant mode.

\paragraph{Ramp-rate limits.}
\begin{equation}
\label{eq:ramp-limits}
  |P_t^{+} - P_{t-1}^{+}| \le R_{\max}\,\Delta t, \qquad
  |P_t^{-} - P_{t-1}^{-}| \le R_{\max}\,\Delta t.
\end{equation}

\paragraph{SOC bounds.}
\begin{equation}
\label{eq:soc-bounds}
  E_{\min} \le E_{t+1} \le E_{\max},
\end{equation}
where $E_{\min} = \mathrm{SOC}_{\min} \cdot E_{\max}$ and similarly for the
upper bound.  Substituting the dynamics~\eqref{eq:soc-dynamics-battery}:
\begin{equation}
\label{eq:soc-bounds-expanded}
  E_{\min} - E_t \;\le\;
    \eta^{+}\,P_t^{+}\,\Delta t - \frac{P_t^{-}}{\eta^{-}}\,\Delta t
  \;\le\; E_{\max} - E_t.
\end{equation}

% ----------------------------------------------------------------
\subsection{Candidate, Safe, and Executed Actions}

Three action vectors are distinguished at each control step:

\begin{description}
  \item[$a_t^\star = (P_t^{+\star}, P_t^{-\star})$:]
    The \emph{candidate action} produced by the dispatch optimiser, maximising
    expected economic value subject to the forecast distribution and the
    \emph{observed} SOC~$\tilde{E}_t$.

  \item[$a_t^{\mathrm{safe}} = (P_t^{+,\mathrm{safe}}, P_t^{-,\mathrm{safe}})$:]
    The \emph{safe action} produced by the DC3S shield
    (Chapter~\ref{ch:dc3s-adapter}).  This is the nearest feasible point to
    $a_t^\star$ within the tightened safe-action set~$\mathcal{A}_t$.

  \item[$a_t^{\mathrm{exec}}$:]
    The \emph{executed action} sent to the BMS.  Under nominal operation
    $a_t^{\mathrm{exec}} = a_t^{\mathrm{safe}}$; under BMS override
    $a_t^{\mathrm{exec}}$ may be further clipped by hardware limits.
\end{description}

The \emph{intervention rate} (IR) measures the fraction of steps where
$a_t^{\mathrm{safe}} \neq a_t^\star$, quantifying the economic cost of
safety enforcement.

% ----------------------------------------------------------------
\subsection{Dispatch Objective}
\label{sec:dispatch-objective}

The deterministic dispatch LP over a rolling horizon~$H$ is:
\begin{equation}
\label{eq:dispatch-lp}
\begin{aligned}
  \max_{P_{t:t+H}^{+},\, P_{t:t+H}^{-}} \quad &
    \sum_{\tau=t}^{t+H-1}
    \bigl[\,\hat{\pi}_\tau\,(P_\tau^{-} - P_\tau^{+})
           - c_{\mathrm{deg}}\,(P_\tau^{+} + P_\tau^{-})\,\bigr]\,\Delta t \\
  \text{s.t.} \quad &
    E_{\tau+1} = E_\tau + \eta^{+}\,P_\tau^{+}\,\Delta t
                       - P_\tau^{-}/\eta^{-}\,\Delta t,
      \quad \forall\,\tau, \\
  & E_{\min} \le E_{\tau+1} \le E_{\max},
      \quad \forall\,\tau, \\
  & 0 \le P_\tau^{+} \le P_{\max}^{+},\;
    0 \le P_\tau^{-} \le P_{\max}^{-},
      \quad \forall\,\tau, \\
  & |P_\tau^{+} - P_{\tau-1}^{+}| \le R_{\max}\,\Delta t,
      \quad \forall\,\tau, \\
  & |P_\tau^{-} - P_{\tau-1}^{-}| \le R_{\max}\,\Delta t,
      \quad \forall\,\tau, \\
  & E_t = \tilde{E}_t \quad (\text{initial condition from observation}).
\end{aligned}
\end{equation}
Here $\hat{\pi}_\tau$ is the forecast price, $c_{\mathrm{deg}}$ is a
degradation penalty per MWh throughput, and the initial SOC~$\tilde{E}_t$ is
the \emph{observed} (possibly degraded) state.

The objective trades revenue (buy low, sell high) against degradation cost.
Only the first action $(P_t^{+\star}, P_t^{-\star})$ is committed; the
remainder is discarded and re-solved at $t+1$ (receding-horizon control).

% ----------------------------------------------------------------
\subsection{Uncertainty and Feasible Set Under Degraded Observation}

When the observation $\tilde{E}_t$ is degraded, the initial-condition
constraint $E_t = \tilde{E}_t$ in~\eqref{eq:dispatch-lp} becomes uncertain.
Let $\hat{E}_t$ denote the controller's best estimate and
$\delta_t = E_t - \hat{E}_t$ the estimation error.  The true SOC after action
is:
\begin{equation}
\label{eq:soc-uncertain}
  E_{t+1}^{\mathrm{true}} = (\hat{E}_t + \delta_t)
    + \eta^{+}\,P_t^{+}\,\Delta t - \frac{P_t^{-}}{\eta^{-}}\,\Delta t.
\end{equation}
For the SOC bounds to hold with probability at least $1 - \alpha$, the
constraint must be tightened by the conformal prediction half-width~$q_t$:
\begin{equation}
\label{eq:tightened-bounds}
  E_{\min} + q_t \;\le\; E_{t+1}^{\mathrm{est}} \;\le\; E_{\max} - q_t,
\end{equation}
where $q_t$ is the reliability-inflated interval half-width from RAC-Cert
(Chapter~\ref{ch:dc3s-adapter}).  This tightening shrinks the feasible set
and reduces the economic value of the dispatch, but guarantees safety
against observation error.

% ----------------------------------------------------------------
\subsection{Calibration Gap and Uncertainty Mismatch}

Even with conformal intervals, a \emph{calibration gap} arises when marginal
coverage differs from conditional coverage across reliability groups.  The
locked CQR group-coverage results
(\texttt{reports/publication/cqr\_group\_coverage.csv}) reveal that wind
intervals achieve only 39--49\% coverage in the low and high reliability
groups despite 90\% nominal marginal coverage.  Solar intervals are better
calibrated (82--90\%) but still exhibit under-coverage at the high-variability
group.

This calibration gap means that the tightened bounds
in~\eqref{eq:tightened-bounds} are \emph{insufficiently tight} during periods
of poor observation quality, motivating the adaptive inflation law in DC3S
(Section~\ref{sec:dispatch-objective} cross-references
Chapter~\ref{ch:dc3s-adapter}).

% ----------------------------------------------------------------
\subsection{Baseline Policies}

The ORIUS evaluation includes eight dispatch controllers spanning four
design philosophies:

\begin{enumerate}
  \item \textbf{Deterministic LP} (\texttt{deterministic\_lp}):
    solves~\eqref{eq:dispatch-lp} on point forecasts with no interval
    awareness.
  \item \textbf{Robust fixed-interval} (\texttt{robust\_fixed\_interval}):
    worst-case robust LP using fixed (non-adaptive) prediction intervals.
  \item \textbf{CVaR interval} (\texttt{cvar\_interval}):
    CVaR-optimised scenario dispatch over quantile forecasts.
  \item \textbf{ACI conformal} (\texttt{aci\_conformal}):
    adaptive conformal inference baseline with online width adjustment.
  \item \textbf{Scenario robust} (\texttt{scenario\_robust}):
    scenario-based robust optimisation over sampled forecast paths.
  \item \textbf{Scenario MPC} (\texttt{scenario\_mpc}):
    receding-horizon model predictive control with scenario tree.
  \item \textbf{DC3S wrapped} (\texttt{dc3s\_wrapped}):
    full DC3S shield with linear reliability inflation.
  \item \textbf{DC3S FTIT} (\texttt{dc3s\_ftit}):
    DC3S with fault-tolerant invariant tube extension.
\end{enumerate}

Results across these controllers are reported in
Chapter~\ref{ch:dc3s-adapter} and the benchmark track
(Chapter~\ref{ch:cpsbench}).

% ----------------------------------------------------------------
\subsection{Scope Boundary}

The dispatch model in this dissertation is a \emph{single-asset, energy-only}
formulation:

\begin{itemize}
  \item \textbf{Included:} energy arbitrage, SOC management, degradation
    penalty, safety constraints.
  \item \textbf{Excluded:} ancillary services (frequency regulation, spinning
    reserve), multi-asset portfolio dispatch, network constraints, and
    market-clearing co-optimisation.
\end{itemize}

The single-asset scope is sufficient to demonstrate the observation-quality
feedback mechanism; extension to fleet-level dispatch is discussed in
Chapter~\ref{ch:orius-bench-battery}.

% ----------------------------------------------------------------
\subsection{Chapter Summary}

This chapter formalised the battery state-space model, the SOC dynamics
equation~\eqref{eq:soc-dynamics-battery}, the power/ramp/SOC constraint set, and the
deterministic dispatch LP~\eqref{eq:dispatch-lp}.  Three action types---candidate,
safe, and executed---were distinguished to precisely locate the safety
intervention point.  The effect of degraded observation on the feasible set
was derived, showing that the conformal half-width must tighten the SOC bounds
and that calibration gaps in wind and solar intervals motivate adaptive
inflation.  Eight baseline controllers spanning deterministic, robust,
conformal, and shielded design philosophies were enumerated.

